\begin{table*}
\centering
\caption{Proposition-level positioning against direct methodological
antecedents.  The middle column identifies what is inherited; the last column
states the remaining boundary of the present contribution.}
\label{tab:literature}
\small
\begin{tabularx}{\textwidth}{p{0.19\textwidth}p{0.34\textwidth}X}
\toprule
Reference & Established antecedent used here & Boundary relative to the present work \\
\midrule
Prud'homme et al. \cite{prudhomme2002reliable} and Quarteroni et al. \cite{quarteroni2011certified}
& Galerkin best approximation, compliant-output error, coercivity-scaled residual bounds, and affine offline--online evaluation
& Not homogenization-specific and not a six-output Fourier residual compiler \\
Boyaval \cite{boyaval2008homogenization}
& Reduced-basis homogenization with output error control and offline--online separation
& Scalar elliptic homogenized outputs rather than the present 3D constitutive transfer realization \\
Yvonnet and He \cite{yvonnet2007r3m}, Fritzen and Leuschner \cite{fritzen2013hybrid}, Hern\'andez et al. \cite{hernandez2014highperformance}, Michel and Suquet \cite{michel2016variational}, and Kunc and Fritzen \cite{kunc2019finite}
& POD/RB, variational reduction, and hyperreduced computational homogenization for constitutive response
& Establish that reduced, structure-preserving constitutive homogenization is not new; they do not provide the present linear-parametric six-output Fourier bound compiler \\
Baur et al. \cite{baur2011interpolatory}
& Parameter value, gradient, and Hessian interpolation under projection-space conditions
& Supplies the classical interpolation principle, not the static six-load constitutive specialization \\
Kami\'nski \cite{kaminski2003sensitivity}, Guo et al. \cite{guo2021constitutive}, and Luke\v{s} and Rohan \cite{lukes2026sensitivity}
& Sensitivity of homogenized properties, reduced constitutive surrogates with derivatives, and sensitivity-guided hyperelastic reduction
& Establish that sensitivities and inverse/design use are not new; they do not provide exact six-load snapshot-containment interpolation with the present bound compiler \\
Huynh et al. \cite{huynh2013scrbe}, Eftang and Patera \cite{eftang2013port}, and Smetana \cite{smetana2015certification}
& Reduced Schur systems, Schur-complement perturbation error, port reduction, and static-condensation certification
& Different component/interface problem; no constitutive FFT two-kernel \(BB/BK/KK\) compilation \\
Huynh et al. \cite{huynh2007scm} and Hain et al. \cite{hain2019hierarchical}
& SCM coercivity lower bounds and hierarchical estimators based on nested reduced approximations
& Classical alternatives/ingredients; the present realization adds an exact matrix Schur drop to a compiled six-output tail bound \\
Kerfriden et al. \cite{kerfriden2014cre} and Hoang et al. \cite{hoang2016twofield}
& CRE and goal-oriented/two-field certification for homogenization or affinely parametrized elasticity
& Different equilibrated two-field estimator family \\
Moulinec and Suquet \cite{moulinec1998numerical} and Vondrejc et al. \cite{vondrejc2014fft,vondrejc2015bounds}
& Isotropic Fourier Green operators, Fourier--Galerkin homogenization, and discretization bounds
& No reduced six-output residual compiler or ROM-error bound of the form proposed here \\
Kochmann et al. \cite{kochmann2019sparsefft}, Vondrejc et al.
\cite{vondrejc2020lowrank}, and Hauck et al. \cite{hauck2026sfft}
& FFT acceleration/reduction by sparse frequency sampling, low-rank tensors,
and tensor-train Fourier transforms
& Reduce the Fourier solve or field representation rather than compiling a
parametric constitutive Schur map and its online matrix bound \\
Pham and Lee \cite{pham2024rbhm}
& Reduced-basis thermal and elastic periodic homogenization with offline--online corrector gradients
& No matrix-valued residual certification or two-kernel dense online estimator \\
Present method
& Uses the classical results above
& Compiles all affine residual pairings for six constitutive outputs into transverse/longitudinal \(BB/BK/KK\) contractions and couples them to physical causal enrichment \\
\bottomrule
\end{tabularx}
\end{table*}
